\documentclass[letter,12pt]{article}

\usepackage[T1]{fontenc}
\usepackage{lmodern}
\usepackage{textcomp}
\renewcommand*\familydefault{\sfdefault}

\usepackage[spanish]{babel}
\usepackage[utf8x]{inputenc}

\usepackage[pdftex]{graphicx}
\usepackage{pifont}
\usepackage[
pdfauthor={Carlos Caballero Burgoa},%
pdftitle={Proyecto Yachay},%
colorlinks,%
citecolor=black,%
filecolor=black,%
linkcolor=black,%
%urlcolor=black
pdftex]{hyperref}

\usepackage{fancyhdr}
\usepackage{lastpage}
\pagestyle{fancy}

% Para la primera página
\fancypagestyle{plain}{
\fancyhead[l]{}
\fancyhead[r]{}
\fancyhead[c]{}
\renewcommand{\headrulewidth}{0.5pt}
\fancyfoot[l]{SCESI \\ Sociedad Científica de Estudiantes de Sistemas e Informática}
\fancyfoot[c]{}
\fancyfoot[r]{\thepage/\pageref{LastPage}}
\renewcommand{\footrulewidth}{0.5pt}}

% Para el resto de páginas
\lhead{Proyecto Yachay}
\chead{}
\rhead{\includegraphics[width=0.1\textwidth]{scesi.png}}
\renewcommand{\headrulewidth}{0.4pt}
\lfoot{SCESI \\ Sociedad Científica de Estudiantes de Sistemas e Informática \\ 
\url {http://scesi.fcyt.umss.edu.bo}}
\cfoot{}
\rfoot{\thepage/\pageref{LastPage}}
\renewcommand{\footrulewidth}{0.4pt}

\title{\bf Yachay}
\author{Carlos Caballero Burgoa} 

\begin{document}
\maketitle
\begin{center}\includegraphics[width=0.48\textwidth]{yachay.png}\end{center}
\begin{center}\url {http://scesi.fcyt.umss.edu.bo}\end{center}
\pagebreak

\tableofcontents
\pagebreak

\section{Introducción}

\section{Antecedentes}


\section{Definición del Problema}
La sociedad cientifica de sistemas e informatica ha venido proveyendo una variedad cada vez mas
amplia de servicios; ya sea desarrollando con fines academicos muchos de ellos, o implantando
soluciones conocidas de software libre. Esta creciente variedad por el momento no se encuentra
integrada en una plataforma común.

Con la llegada de nuevos integrantes y el fomento hacia el uso y desarrollo de software libre, que
es meta de esta sociedad, el problema antes mencionado comezará a ser un problema en lo referente a
interconectividad de estos proyectos.

Una de las barreras para solucionar este problema potencial será el grado de conocimiento que se
debe tener respecto a los proyectos y sus formas de interrelación. Tal barrera sería superada
aprovechando el uso de la plataforma yachay, debido en gran parte a la afinidad y facilidad de
comunicación con los creadores de tal herramienta.

Por lo mencionado anteriormente se define el problema como:

La gran variedad de software potencial utilizado y por utilizar en la sociedad cientifica conducirá 
en breve a un conjunto muy dispar de servicios inconexos.

\section{Objetivo General}
Implantar sobre la infraestructura de la sociedad cientifica, la plataforma yachay para fomentar la
utilización y desarrollo de este proyecto.

\section{Objetivos Específicos}
\begin{itemize}
\item Crear un marco de trabajo en el que puedan ser captadas todas las iniciativas de tipo local.
\item Aprovechar el conocimiento acerca de la herramienta, para proveer funciones especializadas.
\item Unificar todos los servicios que ofrece actualmente y los que ofrecerá en un futuro la
sociedad cientifica.
\end{itemize}

\section{Herramientas}

\section{Justificación}

\begin{thebibliography}{99}
\end{thebibliography}

\end{document}          
